\section{模型评价与敏感性分析}

\subsection{模型优点}

\begin{enumerate}[label=(\arabic*)]
    \item \textbf{强调保真而非单纯降噪。} 第一问通过人工已知信号注入显式检查视觉响应幅值和峰时是否被算法破坏。
    \item \textbf{避免信息泄漏。} 学习型预处理和超参数选择在训练折内完成，并将历史测试结果与开发阶段结果区分。
    \item \textbf{统计结论较克制。} 对共同响应、左右差分、P300样晚期成分分别陈述，不以单一显著性结果代替机制结论。
    \item \textbf{便于向后续问题传递。} 第一问输出统一的试次级EEG、ERP和拟合曲线，可直接作为问题二、三的观测接口。
\end{enumerate}

\subsection{模型局限}

\begin{enumerate}[label=(\arabic*)]
    \item 仅有F3、Fz、F4三个通道，不能完成唯一的三维神经源定位；
    \item 当前幅值缺少可靠物理单位换算，因此只在记录单位下解释；
    \item 左右刺激差分的分半稳定性较弱，需要在问题二利用联合时空动力学而非单点均值进一步建模；
    \item 部分晚期响应对滤波边界和慢漂移处理敏感，因此P300生理解释应保持谨慎。
\end{enumerate}

\subsection{敏感性分析}

\TODO{整合问题一已有的滤波、填充、块长度敏感性；问题二、三完成后补充关键参数扰动与模型结构消融。}
